\cotefile{12}
\scanpage{102}
The association of a cohomology class with a cycle has all the functorial and compatibility properties one could hope for. In particular,if $f:X\handwritten{\longrightarrow}Z$ is a morphism of schemes over $k$,then there are commutative diagrams
\handwritten{\diagram{\begin{tikzcd}[ampersand replacement=\&]
H_\ell^{2i}(Z) \arrow[r,"f^*"] \& H_\ell^{2i}(X) \& H_\ell^{2i}(X) \arrow[r,"f_*"] \& H_\ell^{2i-2r}(Z) \\
C^i(Z) \arrow[u,"\rho_Z"] \arrow[r,"f^*"] \& C^i(X) \arrow[u,"\rho_X"'] \& C^i(X) \arrow[u,"\rho_X"] \arrow[r,"f_*"] \& C^{i-r}(Z) \arrow[u,"\rho_Z"']
\end{tikzcd}}}
where $r=\dim X-\dim Z$.

\uline{Correspondences}. Let $X$ and $Z$ be projective,smooth, and connected schemes over $k$. Then there are canonical isomorphisms
\[
\begin{aligned}
\handwritten{H_\ell^*(X\times Z)}&\handwritten{\xrightarrow{\sim}H_\ell^*(X)\otimes_{\unsure{Q|\mathbb{Q}}_\ell}H_\ell^*(Z)\quad\text{(Kunneth formula)}}\\
&\handwritten{\xrightarrow{\sim}\check{H}_\ell^*(X)\otimes_{Q_\ell}H_\ell^*(Z)\quad\text{(Poincaré duality)}}\\
&\handwritten{\xrightarrow{\sim}\operatorname{Hom}_{Q_\ell}(H_\ell^*(X),H_\ell^*(Z)).}
\end{aligned}
\]
Thus to an element $\handwritten{\alpha}$ of $H_{\handwritten{\ell}}^*(X\mathbin{\mathrm{x}}Z)$ we can canonically associate a homomorphism $\handwritten{\widetilde{\alpha}}:H_{\handwritten{\ell}}^*(X)\handwritten{\longrightarrow}H_{\handwritten{\ell}}^*(Z)$,and to a homomorphism $\handwritten{\beta}:H_{\handwritten{\ell}}^*(X)\handwritten{\longrightarrow}H_{\handwritten{\ell}}^*(Z)$ we can canonically associate an element $\handwritten{\widehat{\beta}\in}H_{\handwritten{\ell}}^*(X\mathbin{\mathrm{x}}Z)$. In particular,we say that $\handwritten{\beta}:H_{\handwritten{\ell}}^*(X)\handwritten{\longrightarrow}H_{\handwritten{\ell}}^*(Z)$ is \uline{algebraic} if the associated cohomology class $\handwritten{\widehat{\beta}}$ is \uline{algebraic},i.e. belongs to $C_{\handwritten{\ell}}^*(X\mathbin{\mathrm{x}}Z)$.

It is sometimes convenient to note that,given $\handwritten{\alpha\in H_\ell^*(X\times Z)}$,the homomorphism $\handwritten{\widetilde{\alpha}}:H_{\handwritten{\ell}}^*(X)\handwritten{\longrightarrow}H_{\handwritten{\ell}}^*(Z)$ is the composition of the homomorphisms
% archive sheet number: 101

\scanpage{117}
\begin{flushright}
Le 27 Mars 1968\\
17H 45\\
\end{flushright}

\begin{center}
\uline{Monsieur GROTHENDIECK}\\
\end{center}

\begin{enumerate}[label=\arabic*.]
\item Melle Zabardi me dit avoir refait dans le détail 6 pages de vos 17 pages de corrections. Elle pense avoir fini demain soir.

\item Mme Hakim ayant appris que les stencils sont stockés Rue de Lille a demandé à Melle Chardon 5 exemplaires de ``La vie Mathématiques en République Démocratique du Vietnam'' et propose de tirer des stencils électroniques.

\item Nous ferons déposer vendredi à Massy les stencils de votre manuscrit et de celui de Jussila.
\end{enumerate}

\begin{flushright}
\handwritten{\unsure{Bonsoir} Schurik.}\\
\handwritten{\uline{\unsure{AR}.}}\\
\end{flushright}
% archive sheet number: 116

\scanpage{121}
une sorte de synthèse géométrico-arithmétique des nombreuses théories cohomologiques dont on dispose\deleted{\unsure{rait}} à présent pour les variétés algébri\deleted{\unsure{qu}}ques,\deleted{de toutes dimensions,}et des représentations linéaires de groupes de \deleted{g}Galois arithmétiques auxquelles ces théories donnent naissance\deleted{\unsure{s}}.

5) \uline{Solution du problème d'introduction d'une} $\handwritten{\lambda}$-\uline{structure} (correspondant à la notion de puissance extérieure) \uline{dans l'anneau} $K^\bullet(X)$(dit ``de Grothendieck'')d'un espace topologique localement annelé (par exemple une variété analytique, ou\deleted{\unsure{nd|md}} \unsure{différentiable}, ou un schéma)\handwritten{,} construit à l'aide de la catégorie des complexes ``parfaits'' sur $X$. Ce problème était signalé comme le plus important problème \inserted{\handwritten{resté}} en suspens dans le Séminaire de Géomé\deleted{y}\handwritten{t}rie Algébrique de l'IHES 1966/67, sur le théorème de Riemann-Roch. La nécessité de travailler avec le $K^\bullet(X)$ défini en termes de \uline{complexes} parfaits, plutôt que de \uline{fibrés} vectoriels seulement, est \unsure{imposé|imposée} par des \deleted{nécess} impératifs techniques irrécusables, qui finiront certainement par s'imposer aussi en dehors de la \deleted{\unsure{G}}\handwritten{g}éométrie algébrique. La solution de Deligne, comme prévu, passe par l'intermédiaire d'une construction de \deleted{\unsure{puis}} foncteurs puissances extérieures dans la catégorie dérivée $D(X)$ (catégorie de\deleted{s} complexes de modules \inserted{\handwritten{sur $X$}}, - en fait, il faut se borner ici aux complexes de \uline{chaines}), en utilisant la théorie de Dold\inserted{\handwritten{-Puppe}}. La difficulté sérieuse consistait à prouver que ces opérations dans $K^\bullet(X)$ satisfont aux propriétés algébriques habituelles, permettant d'exprimer par exemple $\handwritten{\lambda}^i(xy)$ comme polynôme universel à coéfficients entiers en les $\handwritten{\lambda}^j(x),\handwritten{\lambda}^k(y)$. Deligne procède par une voie très ingénieuse, en déterminant d'abord la structure \unsure{da|du} $\handwritten{\lambda}$-anneau de toutes les ``opérations tensorielles universelles'' qu'on peut effectuer sur des fibrés vectoriels. \deleted{\unsure{E}}Il n'a écrit sa théorie que lorsqu'on travaille au dessus d'un corps de base \deleted{\unsure{finis|fini|choisi}}(de caractéristique arbitraire). \deleted{Le cas de la caractéristique}
% archive sheet number: 120

\scanpage{123}
4) \inserted{\handwritten{\uline{\unsure{Trivialite|Trivialité} cohomologique d'}}}\uline{\unsure{Un|un} morphisme projectif et lisse} $f:X\handwritten{\longrightarrow}Y$ \uline{de schémas} (correspondant intuitivement à une fibration sur $Y$ à fibres de variétés projectives non singulières). \deleted{est toujours cohomologiquement trivial (i.e.} \inserted{\handwritten{En d'autres termes,}} \deleted{l}\handwritten{L}a suite spectrale de Leray pour la cohomologie $\handwritten{\unsure{\ell|l}}$-adique - ou à coéffi\deleted{\unsure{s}}cients réels dans le cas du corps de base $C$\deleted{)} - dégénère\deleted{)}. \deleted{Un cas très particulier de ce résultat}\inserted{\handwritten{\deleted{\unsure{E}}Ce théorème est prouvé par Deligne modulo le ``Théorème de \unsure{Lefschetz}'' pour la cohomologie des fibres.}} \handwritten{S}a démonstration s'applique également dans le contexte des variétés analytiques, pour une fibration analytique complexe à fibres compactes kählériennes. Un cas très particulier de ce théorème, relatif au cas où le système local de la cohomologie des fibres était trivial, constitu\deleted{e}\inserted{ait} le résultat principal de la thèse de Blanchard; mais l'hypothèse de Blanchard n'est réalisée dans presque aucun cas intéressant en géométrie. La démonstration de Deligne est différente bien entendu de celle de Blanchard, que Deligne ne connaissait pas.

Déjà du point de vue transcendant, l'intérêt du résultat de Deligne \deleted{\unsure{en top}} pour la topologie des variétés analytiques est évident. Dans l'optique du rédacteur cependant, ce sont ses implications en géométrie algébrique abstraite,et à longue échéance, à la théorie des nombres, qui sont les plus intéressantes\handwritten{,} \deleted{Par exemple, le}\deleted{\unsure{et}} \handwritten{C}e sont d'ailleurs des considérations arithmétiques qui avaient amené le \deleted{rédacteur}\inserted{\handwritten{rapporteur}} à communiquer à Deligne \deleted{la conjecture} l'énoncé précédent \unsure{somme|comme} une conjecture. Grâce au théorème de Deligne, on peut voir par exemple que les hypothèses de \unsure{Reimann|Riemann}-Weil en géométrie diophantienne, interprétées en termes des propriétés d'un endomorphisme de frobénius sur la cohomologie $\handwritten{\ell}$-adique, \deleted{restent valables sous une} impliquent des propriétés analogues dans le cas de coéfficients tordus (pouvant avoir des singularités quelconques, de plus). Le théorème de Deligne sera un des \inserted{\handwritten{nombreux}} ingrédients techniques \deleted{nécessaires}\inserted{\handwritten{indispensables}} pour développer la ``théorie des motifs'', qui \handwritten{\unsure{sera}}
% archive sheet number: 122

\scanpage{126}
\mspage{-III- \handwritten{7} -}
\begin{flushright}n° \handwritten{383}\end{flushright}
quotient un groupe essentiel dans l'espace euclidien $E/W'$.

On dit que $G$ est \uline{réductible} si $E$ est somme directe de deux sous-espaces orthogonaux $E',E''$ distincts de $E$ et stables pour ${}^vG$. Alors tout élément $s$ de $G$ définit, par passage au quotient, des automorphismes affines $s'$ de $E/E''$ (identifié à $E'$) et $s''$ de $E/E'$ (identifié à $E''$) ; l'ensemble des $s'$ (resp. $s''$) est un sous-groupe $G'$ (resp. $G''$) du groupe affine de $E'$ (resp. $E''$). Si $H$ est un hyperplan de $E$ tel que $s_H\in G$, on voit aisément que, ou bien $E'$ est contenu dans la direction ${}^vH$ de $H$, ou bien $E'$ contient la droite orthogonale à $H$ ; dans le premier cas l'automorphisme $s'$ de $E'$ défini par $s_H$ est l'identité, et dans le second $s'$ est la réflexion par rapport à l'hyperplan $H\handwritten{\cap}E'$ de $E'$. Il s'ensuit que $G'$ et $G''$ sont engendrés par des réflexions. D'autre part, l'application $s\handwritten{\rightsquigarrow}(s',s'')$ de $G$ dans $G'\handwritten{\times}G''$ est évidemment un isomorphisme.

On dit que $G$ est \uline{irréductible} s'il n'est pas réductible. Notant $\underline{H}$ l'ensemble des hyperplans tels que $s_H\in G$, ceci revient à dire qu'il n'existe pas de partition de $\underline{H}$ en deux sous-ensembles tels que tout hyperplan du premier soit perpendiculaire à tout hyperplan du second. Dans le cas général il résulte de ce qu'on vient de voir que $E$ s'identifie khanoniquement à un produit d'espaces euclidiens $E^{(1)},\ldots,E^{(q)}$, et $G$ à un produit de groupes $G^{(1)},\ldots,G^{(q)}$ opérant dans ces espaces, de telle sorte que chaque $G^{(i)}$ soit irréductible.

Nous nous fixerons désormais un ensemble $\underline{H}$ d'hyperplans (affines) de $E$ tel que

1) les réflexions $s_H$ $(H\in\underline{H})$ engendrent $G$,

2) les hyperplans de $\underline{H}$ sont permutés par $G$.

Comme au n°1, notons que, pour $g\in G$, on a $s_{g(H)}=gs_Hg^{-1}$. Nous verrons plus loin que l'ensemble des $s_H$ $(H\in\underline{H})$ est l'ensemble de toutes les réflexions de $G$.

PROPOSITION 4.- \uline{La famille d'hyperplans} $\underline{H}$ \uline{est localement finie}.

Raisonnons par l'absurde. Soit $x\in E$ tel que tout voisinage $U$ de $x$ rencontre une infinité d'hyperplans de $\underline{H}$. Alors $g(U)\handwritten{\cap}U\handwritten{\ne\varnothing}$ pour une infinité de réflexions affines
% archive sheet number: 125

\scanpage{128}
\mspage{-III- \handwritten{6} -}
\begin{flushright}n°\handwritten{383}\end{flushright}
COROLLAIRE 7.- \uline{Soient} $C$ \uline{une chambre et} $x\in E$. \uline{Il existe un élément} $y$ \uline{de} $\overline{C}$ \uline{et un seul et un élément} $g\in G$ \uline{tels que} $x=g(y)$.

En effet la réunion des chambres est le complémentaire de $\displaystyle\handwritten{\bigcup}_{H\in\underline{H}}H$, donc est dense dans $E$. Ainsi $x$ est adhérent à une chambre $g(C)$ $(g\in G)$, soit $x=g(y)$ avec $y\in\overline{C}$. L'unicité de $y$ résulte alors du cor.5 au th.2, \S\ 2, n°3.

COROLLAIRE 8.- \uline{Soient} $x\in E$ \uline{et} $G'$ \uline{le sous-groupe de} $G$ \uline{formé des} $g\in G$ \uline{tels que} $g(x)=x$. \uline{Alors} $G'$ \uline{est engendré par les réflexions} $s_H$ $(H\in\underline{H})$ \uline{telles que} $s_H(x)=x$.

On peut supposer $G$ essentiel, et que $x$ est adhérent à la chambre $C$. Alors $x$ appartient à une face $C_J$ $(J\handwritten{\subset\{}1,2,\ldots,m\handwritten{\}})$ de $C$ (avec les notations du \S\ 2, n°3). Pour $g\in G'$, on a $C_J\handwritten{\cap}g(C_J)\handwritten{\ne\varnothing}$, d'où $gG_J=G_J$ (\S\ 2, n°3, th.2), et $G'\handwritten{\subset}G_J$. Comme $G_J$ laisse $C_J$, et donc $x$, invariant, on a $G'=G_J$. Donc $G'$ est engendré par les réflexions $s_i$ pour $i\handwritten{\notin}J$.

2. \uline{Groupes affines engendrés par des réflexions}.

\uline{Dans ce n°}, $G$ \uline{désigne un sous-groupe du groupe affine de} $E$ (Alg.II, 2ème éd., App.II, n°4), \uline{engendré par des réflexions par rapport à des hyperplans affines}. \uline{On suppose que} $G$, \uline{muni de la topologie discrète}, \uline{opère proprement dans} $E$ (Top.Gén.,III, \S\ 4).

A toute application affine $s$, faisons correspondre l'application linéaire ${}^vs$ associée à $s$ (Alg.II, 2ème éd., App.II, n°4). Si $s$ est la réflexion par rapport à un hyperplan $H$, ${}^vs$ est la réflexion (vectorielle) par rapport à la direction ${}^vH$ de $H$ (ibid., n°3). Ainsi $s\handwritten{\rightsquigarrow}{}^vs$ est un homomorphisme de $G$ sur un sous-groupe ${}^vG$ du groupe orthogonal de $E$, engendré par des réflexions.

Soit $W$ le sous-espace essentiel de $E$ pour ${}^vG$, et soit $W'$ l'orthogonal de $W$ dans $E$. Pour tout $x\in E$ et tout $s\in G$, on a $s(x)-x\in W$. Si $W=E$, autrement dit si ${}^vG$ est essentiel, on dit que $G$ est \uline{essentiel}. Dans le cas général, $G$ définit par passage au
% archive sheet number: 127

\scanpage{130}
\mspage{-III- \handwritten{5} -}
\begin{flushright}n°\handwritten{383}\end{flushright}
On a vu en effet que $C$ (ou une de ses transformées) est l'ensemble des $\displaystyle\handwritten{\sum}_{i\in I}a_i e'_i$ tels que $a_i>0$, et que $\cos(\handwritten{\widehat}{e'_i,e'_j})=\cos(\handwritten{\pi}/m_{ij})\handwritten{\geq}0$ (prop.2).

COROLLAIRE 2.- \uline{Pour tout} $H\in\underline{H}$, \uline{soit} $e_H$ \uline{le vecteur unitaire orthogonal à} $H$ \uline{et situé du même côté de} $H$ \uline{que} $C$. \uline{Alors} $e_H$ \uline{est combinaison linéaire à coefficients positifs de} $e_1,\ldots,e_m$.

En effet, d'après la prop.3, on a $\displaystyle e_H=\handwritten{\sum}_{j\in I}b_j e_j$ avec $b_j\in R$, d'où $b_j=(e_H\handwritten{|}e'_j)$. Comme $e'_j\in\overline{C}$, et que $e_H$ est du même côté de $H$ que $C$, on a $(e_H\handwritten{|}e'_j)\handwritten{\geq}0$.

COROLLAIRE 3.- \uline{Avec les mêmes notations}, $s_i$ \uline{laisse stable l'ensemble des} $e_H$ \uline{distincts de} $e_i$.

Pour $e_H\handwritten{\ne}e_i$, on a $\displaystyle e_H=\handwritten{\sum}_{j=1}^m a_j e_j$ où les $a_j$ sont des nombres réels positifs (cor.2), et où l'un des $a_j$, pour $j\handwritten{\ne}i$, est $>0$. Le vecteur $s_i(e_H)=e_H-2(e_H\handwritten{|}e_i)e_i$ est, soit un des $e_{H'}$, soit un de leurs opposés ; or, pour $j\handwritten{\ne}i$, le coefficient de $e_j$ dans la décomposition de $s_i(e_H)$ est $a_j$ ; donc $s_i(e_H)$ ne peut être de la forme $-e_{H'}$ car ces $a_j$ ne sont pas tous $\handwritten{\leq}0$ (cor.2). Ainsi $s_i(e_H)$ est de la forme $e_{H'}$. C.Q.F.D.

COROLLAIRE 4.- \uline{L'ordre de} $G$ \uline{est égal au nombre des chambres définies par} $\underline{H}$.

Ceci résulte aussitôt de b).

COROLLAIRE 5.- \uline{Avec les notations du th.1}, \uline{l'ensemble des réflexions} $s_H$ $(H\in\underline{H})$ \uline{est nécessairement l'ensemble de toutes les réflexions de} $G$.

Soit $\underline{H}'\handwritten{\supset}\underline{H}$ l'ensemble des hyperplans $H$ tels que $s_H\in G$. Toute chambre relative à $\underline{H}'$ est contenue dans une chambre relative à $\underline{H}$ et une seule. D'après le cor.4, les chambres relatives à $\underline{H}'$ et à $\underline{H}$ sont donc les mêmes. D'après le lemme 3 du \S\ 1, n°1, tout élément de $\underline{H}'$ est mur d'une chambre, donc appartient à $\underline{H}$.

COROLLAIRE 6.- \uline{Soient} $C$ \uline{et} $C'$ \uline{deux chambres}. \uline{Alors} $\overline{C}\handwritten{\cap}\overline{C'}$ \uline{est l'adhérence d'une face} $F$ \uline{de} $C$. \uline{L'élément} $g\in G$ \uline{tel que} $C'=g(C)$ \uline{est un produit de réflexions} $s_H$ $(H\in\underline{H})$ \uline{par rapport à des hyperplans} $H$ \uline{contenant} $F$.

C'est une traduction du lemme 7, \S\ 2, n°4.
% archive sheet number: 129

\scanpage{132}
\mspage{-III- \handwritten{4} -}
\begin{flushright}n° \handwritten{383}\end{flushright}
$s_i$, donc pour $G$ (prop.1), et est par conséquent orthogonal à $W$ ; ceci montre que les $e_i$ engendrent $W$. C.Q.F.D.

Ceci étant, soient $I=\handwritten{\{}1,2,\ldots,m\handwritten{\}}$, $M=(m_{ij})$ la matrice de Coxeter sur $I$ où $m_{ij}$ est l'ordre de $s_i s_j$ dans $G$, et $G(M)$ le groupe de Coxeter défini par les générateurs $r_i$ $(i\in I)$ et les relations $(r_i r_j)^{m_{ij}}=1$ (\S\ 2, n°1). D'après la prop.3 on peut prendre $W$, muni de la base $(e_i)$, pour espace de la représentation directe $\handwritten{\rho}$ de $G(M)$ ; d'après la prop.2 et le fait que les $e_i$ sont unitaires, la forme bilinéaire $B_M$ du \S\ 2, n°2 est la restriction à $W$ de la forme donnée $(x\handwritten{|}y)$ ; ainsi $\handwritten{\rho}(r_i)$ est la restriction $s'_i$ à $W$ de la réflexion $s_i$. Identifions $W$ à son dual au moyen de la forme $(x\handwritten{|}y)$ ; comme, pour passer d'une symétrie à sa contragrédiente, on échange vecteur et forme linéaire (\S\ 2, n°3), et comme $s'_i$ est définie par un vecteur unitaire et l'hyperplan orthogonal, $s'_i$ est sa propre contragrédiente. Si $(e'_i)$ est la base de $W$ duale de $(e_i)$ (i.e. $(e_i\handwritten{|}e'_j)=\handwritten{\delta}_{ij}$ pour tout $i,j$), l'ensemble des $\displaystyle x=\handwritten{\sum}_{i\in I}a_i e'_i$ $(a_i\in R)$ de $W$ tels que $(x\handwritten{|}e_i)>0$ pour tout $i$ est celui des $x$ tels que $a_i>0$ pour tout $i$ ; donc, d'après (2), $C\handwritten{\cap}W$ est la chambre notée $C$ au \S\ 2, n°3. En tenant compte du fait que $G$ opère trivialement sur le sous-espace $W'$ orthogonal à $W$, le th.1 du \S\ 2, n°3 donne le résultat suivant :

THEOREME 1.- \uline{Soient} $\underline{H}$ \uline{un ensemble d'hyperplans de} $E$ \uline{vérifiant} (1) \uline{et} (2), $C$ \uline{une chambre définie par} $\underline{H}$, $H_1,\ldots,H_m$ \uline{ses murs}, \uline{et} $s_i$ \uline{la réflexion par rapport à} $H_i$. \uline{Notons} $m_{ij}$ \uline{l'ordre de} $s_i s_j$ \uline{dans} $G$. \uline{Alors}

a) \uline{les hyperplans} $H_i$ \uline{sont linéairement indépendants} ;

b) \uline{le groupe} $G$ \uline{est défini par les générateurs} $s_i$ \uline{et les relations} $s_i^2=(s_i s_j)^{m_{ij}}=1$;

c) \uline{le groupe} $G$ \uline{opère de façon simplement transitive sur l'ensemble des chambres définies par} $\underline{H}$.

COROLLAIRE 1.- \uline{Si} $G$ \uline{est essentiel}, \uline{deux vecteurs quelconques} $x,y$ \uline{d'une même chambre} $C$ \uline{font un angle aigu} (i.e. $(x\handwritten{|}y)\handwritten{\geq}0$).
% archive sheet number: 131

\scanpage{134}
\mspage{-III- \handwritten{3} -}
\begin{flushright}n°\handwritten{383}\end{flushright}
pour $G^{\unsure{o|0|\circ}}$ est le plan $P$ engendré par $e_i$ et $e_j$. Les hyperplans de la forme $g(H_i)$ et $g(H_j)$, pour $g\in G^o$, contiennent tous le sous-espace $P'$ orthogonal à $P$, et forment une partie $\underline{H}^o$ de $\underline{H}$ qui jouit, vis à vis de $G^o$, des propriétés 1) et 2) ci-dessus. Comme $\underline{H}^o\handwritten{\subset}\underline{H}$, il existe une chambre $C^o$ définie par $\underline{H}^o$ et contenant $C$ ; il est clair que $C^o$ est un des dièdres ayant $H_i$ et $H_j$ pour murs. Ainsi $C^o\handwritten{\cap}P$ est un secteur angulaire, limité par deux demi-droites $D_i\handwritten{\subset}H_i$ et $D_j\handwritten{\subset}H_j$. Posant $\handwritten{\unsure{\vartheta|\theta}}=(\handwritten{\widehat}{D_i,D_j})$, $s_j s_i$ est la rotation de $2\handwritten{\vartheta}$ autour de $P'$, et les transformées de $C^o$ par $G^o$ sont ses transformées par les rotations d'angle $n\handwritten{\vartheta}$ $(n\in Z)$ ; comme deux transformées distinctes de $C^o$ sont disjointes, il résulte du \S\ 1, n°3, lemme 9 que $\handwritten{\vartheta}$ est de la forme $\handwritten{\pi}/m$ $(m\in Z)$ ; on a nécessairement $m=m_{ij}$. Par définition de $e_i$ et $e_j$, et avec une orientation convenable de $P$, on a $(\handwritten{\widehat}{D_i,e_i})=-\handwritten{\pi}/2$, et $(\handwritten{\widehat}{D_j,e_j})=+\handwritten{\pi}/2$. D'où, par la formule de Chasles (Alg.IX, \S\ 10, n°3, (9)), $(\handwritten{\widehat}{e_i,e_j})=(\handwritten{\widehat}{e_i,D_i})+(\handwritten{\widehat}{D_i,D_j})+(\handwritten{\widehat}{D_j,e_j})=\handwritten{\pi}+(\handwritten{\widehat}{D_i,D_j})$. On en déduit $(e_i\handwritten{|}e_j)=\cos(\handwritten{\widehat}{e_i,e_j})=-\cos(\handwritten{\widehat}{D_i,D_j})=-\cos(\handwritten{\pi}/m_{ij})$. La seconde assertion résulte de la monotonie de la fonction cosinus dans $[0,\handwritten{\pi}/2]$ et de la formule $\cos(\handwritten{\pi}/3)=1/2$ (référence introuvable !!).

PROPOSITION 3.- \uline{Les vecteurs} $e_1,\ldots,e_m$ \uline{forment une base du sous-espace essentiel} $W$ \uline{de} $E$.

En effet l'ensemble $C'$ formé par $0$ et par les $y\in E$ tels que $(y\handwritten{|}x)>0$ pour tout $x\in C$ est un cône convexe pointé, et on a $C'\handwritten{\cap}(-C')=\handwritten{\{}0\handwritten{\}}$ ; donc $C'$ est l'ensemble des éléments positifs pour une relation d'ordre faisant de $E$ un espace vectoriel ordonné (E.V.T., II, \S\ 1, n°5). D'après (2) on a $e_i>0$ pour tout $i$, et d'après la prop.2, on a $(e_i\handwritten{|}e_j)\handwritten{\leq}0$ pour $i\handwritten{\ne}j$. Le lemme 6 du \S\ 1, n°2 montre alors que les $e_i$ sont linéairement indépendants. D'autre part il est clair que les $e_i$ appartiennent à $W$. Enfin si un élément de $E$ est orthogonal à $e_1,\ldots,e_m$, il est fixe pour les
% archive sheet number: 133

\scanpage{136}
\mspage{-III- \handwritten{2} -}
\begin{flushright}n°\handwritten{383}\end{flushright}
Considérons les \uline{chambres} définies par $\underline{H}$, c'est-à-dire les composantes connexes de la réunion des hyperplans $H\in\underline{H}$ (\S\ 1, n°1) ; rappelons qu'une chambre est un cône ouvert convexe (ibid.).

PROPOSITION 1.- \uline{Soient} $C$ \uline{une chambre}, $H_1,\ldots,H_m$ \uline{ses murs}. \uline{Les réflexions} $s_{H_1},\ldots,s_{H_m}$ \uline{engendrent} $G$. \uline{Le groupe} $G$ \uline{opère transitivement sur l'ensemble des chambres définies par} $\underline{H}$.

Soient $G'$ le sous-groupe de $G$ engendré par $s_{H_1},\ldots,s_{H_m}$, $x$ un point de $C$, $C'$ une autre chambre, et $x'$ un point de $C'$. Parmi tous les transformés $g'(x')$ $(g'\in G')$, soit $y$ l'un de ceux qui sont les plus voisins de $x$. Si $y$ n'était pas dans $C$, il existerait un mur $H_i$ de $C$ tel que $x$ et $y$ soient de part et d'autre de $H_i$ ; alors le symétrique $s_{H_i}(y)$ serait strictement plus proche de $x$ que $y$ ; on a donc $y\in C$, ce qui montre que $G'$ opère transitivement sur l'ensemble des chambres définies par $\underline{H}$, et prouve la seconde assertion.

Soit maintenant $H\in\underline{H}$. D'après le \S\ 1, n°1, lemme 3, $H$ est un mur d'une chambre $C'$ définie par $\underline{H}$. Il existe donc, d'après ce qu'on vient de voir, un élément $g'$ de $G'$ et un indice $i$ $(1\handwritten{\leq}i\handwritten{\leq}m)$ tels que $H=g'(H_i)$. D'après (1) on a donc $s_H=g's_{H_i}{g'}^{-1}$, d'où $s_H\in G'$ ; ceci montre que $G=G'$. C.Q.F.D.

Nous fixerons désormais une chambre $C$ (définie par $\underline{H}$), noterons $H_1,\ldots,H_m$ ses murs, et poserons $s_{H_i}=s_i$. Soit $e_i$ le vecteur unitaire orthogonal à $H_i$ et situé du même côté de $H_i$ que $C$. La chambre $C$ est donc l'ensemble des points $x$ de $E$ tels que
\[
(x\handwritten{|}e_i)>0\qquad\text{pour tout }i=1,\ldots,m.
\tag*{(2)}
\]

PROPOSITION 2.- \uline{Soient} $i,j$ \uline{deux indices distincts}, \uline{et} $m_{ij}$ \uline{l'ordre de} $s_i s_j$ \uline{dans} $G$. \uline{On a alors} $(e_i\handwritten{|}e_j)=-\cos(\handwritten{\pi}/m_{ij})$, \uline{de sorte que} $(e_i\handwritten{|}e_j)\handwritten{\leq}0$. \uline{Si} $(e_i\handwritten{|}e_j)<0$, \uline{on a} $(e_i\handwritten{|}e_j)\handwritten{\leq}-1/2$.

Soit $G^{\unsure{o|0|\circ}}$ le sous-groupe de $G$ engendré par $s_i$ et $s_j$ ; le sous-espace essentiel de $E$
% archive sheet number: 135

\scanpage{138}
\mspage{-III- 1 -}
\begin{flushright}n° \handwritten{383}\end{flushright}
\begin{center}\S\ 3. \uuline{Groupes engendrés par des réflexions}.\end{center}

\uline{Dans ce} \S, \uline{on désigne par} $E$ \uline{un espace hilbertien réel de dimension finie} $n$, \uline{par} $(x\handwritten{|}y)$ \uline{sa forme bilinéaire}, \uline{et par} $\handwritten{\lVert}x\handwritten{\rVert}$ \uline{sa norme}. \uline{Pour tout hyperplan} (vectoriel ou affine) $H$ \uline{de} $E$, \uline{on note} $s_H$ \uline{la réflexion} (\uline{ou symétrie orthogonale}) \uline{par rapport à} $H$ (cf. \S\ 2, n°2).

1. \uline{Groupes finis engendrés par des réflexions vectorielles}.

\uline{Dans ce} n°, $G$ \uline{désigne un sous-groupe fini du groupe orthogonal de} $E$, \uline{engendré par des réflexions} (par rapport à des hyperplans vectoriels).

Soit $\underline{H}$ un ensemble fini d'hyperplans de $E$ tel que les $s_H$ $(H\in\underline{H})$ engendrent $G$. L'ensemble $W'$ des $x\in E$ invariants pour $G$ est aussi l'ensemble des $x\in E$ invariants par les $s_H$ $(H\in\underline{H})$ ; donc $\displaystyle W'=\handwritten{\bigcap}_{H\in\underline{H}}H$. Le sous-espace $W$ de $E$ orthogonal à $W'$ est stable pour $G$ ; on l'appelle le sous-espace \uline{essentiel} de $E$ (pour $G$). L'application qui à tout $g\in G$ fait correspondre sa restriction à $W$ est un isomorphisme de $G$ sur un sous-groupe fini $G'$ du groupe orthogonal de $W$, engendré par les réflexions $s_{H\handwritten{\cap}W}$ $(H\in\underline{H})$. On dit que $G$ est \uline{essentiel} si $W=E$ ; en tous cas $G'$ est essentiel, ce qui permet de se ramener le plus souvent au cas où $G$ est essentiel.

Nous nous fixerons désormais un ensemble fini $\underline{H}$ d'hyperplans de $E$ tel que

1) les réflexions $s_H$ $(H\in\underline{H})$ engendrent $G$,

2) les hyperplans de $\underline{H}$ sont permutés par $G$.

Notons, à ce propos, que, pour $g\in G$, on a
\[
s_{g(H)}=gs_Hg^{-1}.
\tag*{(1)}
\]
En principe, toutes les constructions qui suivent dépendent du choix de $\underline{H}$ ; en fait, nous nous abstiendrons de le rappeler dans nos notations et nos énoncés, car le cor.5 du th.1 révélera que l'ensemble des $s_H$ est nécessairement l'ensemble de \uline{toutes} les réflexions de $G$.
% archive sheet number: 137
